\section{Mempool}
\label{sec:mempool}

The mempool is local, non-consensus state. It is a node's current set of
unconfirmed transactions that the node is willing to keep, relay, and consider
for block-template construction. A transaction can be valid in a block and
absent from a node's mempool.

Let \(U\) be the active UTXO set. A mempool is:

\[
  Q = (M,G,I),
\]

where \(M\) is a finite map from transaction identifier to mempool entry, \(G\)
is the dependency graph induced by unconfirmed spends, and \(I\) is the set of
indexes needed to query entries by transaction identifier, witness transaction
identifier, spent outpoint, created outpoint, feerate, and entry time.
In equations, \(t \in M\) ranges over transactions in the entry map, and
\(Q \cup \{tx\}\) denotes the mempool obtained by adding \(tx\) to \(M\) and
recomputing \(G\) and \(I\).
For validation against mempool ancestors, define:

\[
  \operatorname{View}(Q,U)=U \cup \bigcup_{t \in M}\operatorname{Out}(t).
\]

\subsection{Entry}

\begin{longtable}{@{}L{0.18\linewidth}L{0.22\linewidth}L{0.48\linewidth}@{}}
\toprule
Field & Type & Meaning \\
\midrule
\code{tx} & \code{Transaction} & Serialized transaction accepted into the mempool. \\
\code{txid}, \code{wtxid} & \code{uint256} & Transaction identifiers from Section~\ref{sec:transactions}. \\
\code{fee} & \code{int64} & Input value minus output value under \(\operatorname{View}(Q,U)\). \\
\code{weight}, \code{vsize} & \code{int64} & Weight and virtual size from Section~\ref{sec:serialization}. \\
\code{sigopsCost} & \code{int64} & Signature-operation cost under the selected script rules. \\
\code{spends} & OutPoint set & Previous outputs consumed by \code{tx}. \\
\code{creates} & OutPoint set & Outputs created by \code{tx}. \\
\code{parents} & Txid set & Mempool transactions directly spent by \code{tx}. \\
\code{children} & Txid set & Mempool transactions that directly spend \code{tx}. \\
\code{entryTime}, \code{entryHeight} & local clock, height & Admission time and active-chain height at admission. \\
\bottomrule
\end{longtable}

The entry fields are logical fields: an implementation may derive, cache, or
index them differently, but externally observable mempool behavior must be
equivalent.

\subsection{Graph Invariants}

For a transaction \(t\), let \(\operatorname{In}(t)\) be its spent outpoints and
\(\operatorname{Out}(t)\) be its created outpoints. Direct mempool dependencies
are:

\[
\begin{aligned}
\operatorname{parents}_{Q}(t)
  &= \{u \in M : \operatorname{In}(t) \cap \operatorname{Out}(u) \ne \emptyset\},\\
\operatorname{children}_{Q}(t)
  &= \{u \in M : \operatorname{In}(u) \cap \operatorname{Out}(t) \ne \emptyset\}.
\end{aligned}
\]

A mempool is well-formed relative to \(U\) iff:

\[
\begin{aligned}
\operatorname{TxAvailable}_{Q}(t,U) \Longleftrightarrow\;&
  \forall x \in \operatorname{In}(t):
  x \in U \cup \bigcup_{u \in M \setminus \{t\}}\operatorname{Out}(u).
\end{aligned}
\]

\[
\begin{aligned}
\operatorname{WellFormed}(Q,U) \Longleftrightarrow\;&
  \forall t \in M:\neg\operatorname{Coinbase}(t) \wedge
  \operatorname{ContextFreeValid}(t) \wedge
  \operatorname{TxAvailable}_{Q}(t,U) \wedge \\
& \forall a,b \in M,\ a\ne b:
  \operatorname{In}(a) \cap \operatorname{In}(b)=\emptyset \wedge \\
& G=\{(u,t):u \in \operatorname{parents}_{Q}(t)\} \wedge
  \operatorname{Acyclic}(G).
\end{aligned}
\]

The graph order is the spend order: a parent precedes every child. Any block
template or package derived from \(Q\) must include transactions in a topological
order of \(G\).

\subsection{Admission}

Admission is a local transition:

\[
  Q \xrightarrow[P]{tx,U,H} Q',
\]

where \(P\) is the node's local policy parameter set and \(H\) is the next-block
height and median-time context. Write \(H=(C,h,t)\), where \(C\) is the active
chain, \(h\) is the next block height, and \(t\) is the next block locktime
cutoff. Let \(V_Q=\operatorname{View}(Q,U)\). A single transaction admission
succeeds iff:

\[
\begin{aligned}
\operatorname{Admit}_{P}(Q,U,H,tx) \Longleftrightarrow\;&
  tx \notin M \wedge
  \operatorname{WellFormed}(Q \cup \{tx\},U) \wedge \\
& \operatorname{ConsensusValid}(V_Q,tx,H) \wedge
  \operatorname{AbsFinal}(tx,h,t) \wedge
  \operatorname{SeqFinal}(tx,C,V_Q,h) \wedge \\
& \operatorname{Policy}_{P}(Q,U,tx,H),\\
Q' = Q \cup \{tx\} \Longleftrightarrow\;&
  \operatorname{Admit}_{P}(Q,U,H,tx).
\end{aligned}
\]

\(\operatorname{Policy}_{P}\) includes node-local standardness, feerate,
resource, replacement, package, and denial-of-service rules. Those rules do not
define consensus validity; when they are specified, they belong with the
transaction, script, package, or relay behavior they constrain.

A package \(\pi\) is admitted atomically by applying the same transition to a
topologically sorted sequence of transactions. If any member fails, the package
transition fails and \(Q\) is unchanged.

\subsection{Removal}

Let \(B\) be a connected block and \(U'\) the UTXO set after connecting it.
Let \(H'\) be the next-block context after connecting \(B\).
Connecting \(B\) removes every included transaction and every mempool
transaction that conflicts with the block:

\[
\begin{aligned}
\operatorname{ConnectBlock}(Q,B,U') =
\{t \in M:\;&t \notin B.\text{\code{txs}} \wedge \\
&\operatorname{In}(t) \cap \operatorname{Spent}(B)=\emptyset \wedge \\
&\operatorname{ConsensusValid}(\operatorname{View}(Q,U'),t,H')\}.
\end{aligned}
\]

After removal, all entry indexes and dependency edges are recomputed so that
\(\operatorname{WellFormed}(Q,U')\) holds.

Disconnecting a block does not automatically restore its transactions to the
mempool. The disconnected transactions become admission candidates and may
reenter only through the current admission relation under the post-disconnect
chain state.

\subsection{Orphans and Relay View}

A transaction with unavailable inputs is not a mempool entry. A node may keep an
orphan cache \(O\) of such transactions and retry admission when parents become
available, but \(O\) is local resource state and has no consensus meaning.

Transaction relay announces transactions from \(Q\) by \code{txid} or
\code{wtxid}, serves requested transactions if still present, and may filter,
delay, or suppress announcements according to local policy. Relay choices do
not change the definition of \(Q\) \cite{bitcoin-core-v31}.
